% Section 1. Introduction.

Behavioral researchers now prompt large language models (LLMs) with demographic personas and treat the answers as \emph{silicon samples}.
\citet{argyle2023outofone} used such samples to reproduce vote choice and opinion by group, \citet{park2024diminished} used them to replicate classic experiments, and \citet{lin2026simulators} describe their use in piloting instruments before fielding them.
For a study, a simulated sample costs little, reaches groups that are hard to recruit, and can be rerun whenever the design changes.
Each of these uses assumes that the simulated sample stands in for a real one in the respect the study measures.

In practice the stand-in assumption is checked unevenly.
Across eight social science journals, \citet{desai2026validating} found validation practice in LLM studies ``inconsistent and limited.''
\citet{park2024diminished} found GPT-3.5 answering with almost no variation where people vary, and \citet{bisbee2024synthetic} found synthetic opinion distributions narrower than the survey distributions they imitate.
\citet{lin2026workflow} proposed a workflow in which the validity a study needs rises with the claim it makes.
That workflow does not yet supply the tests themselves, each with a pass rule and a population reference, in an order a researcher can run on one model and report.

One persona from the corpus analysed here shows what reading individual answers misses.
Prompted as a single Black woman on Medicaid with a household income under \$35,000, GPT-4o-mini returns a PHQ-8 total of 8 on its first draw, a plausible answer in the mild band. % R: 85_worked_case.csv
Over thirty draws the same prompt averages 10.3. % R: 85_worked_case.csv
Adults in the National Health and Nutrition Examination Survey (NHANES) who share her race, sex, income and marital status average 4.7. % R: 85_worked_case.csv
None of her single answers would show that the simulated population she belongs to sits well above the real one.

We propose a protocol for the missing check, the validity ladder (Figure~\ref{fig:ladder}).
A gate first measures how many repeated draws make a persona's score precise enough to read.
Above the gate, level 1 tests whether a single draw is an answer pattern a person could give, and level 2 tests whether gaps between groups match the population's.
Level 3 tests whether the population level matches, and level 4 tests whether the items relate to one another as they do in people.
Each rung of the ladder, the gate or one of the levels, carries a statistic, a pass rule and the use a pass supports, and a simulated sample supports only the claims whose rungs it has passed.
With the rules frozen and released as code, a researcher can run the ladder on their own sample and report the verdicts.

The worked example applies the ladder to 28,800 PHQ-8 depression assessments from GPT-4o-mini, Gemini-3-Flash, DeepSeek-V3 and GLM-4.7, a corpus first collected for an earlier audit \citep{keough2026plausible} and rebuilt here from the raw answers. % R: corpus_v3_provenance.csv
A few repeated draws suffice to read each persona's mean, yet no model passes any level above the gate.
Single answers vary less than people's, and the models exaggerate the income gradient, place every group's level too high and distort the scale's structure.
Before these failures can be credited to the models, the ladder has to be shown not to fail real people.
Survey respondents drawn through the same design rarely fail any level, and failures planted at known sizes are caught by the level built to find them.

Three published silicon-sample datasets test whether the ladder travels beyond one instrument.
Feeling thermometers from \citet{bisbee2024synthetic}, vote choice and attitudes from \citet{argyle2023outofone} and opinion distributions from \citet{meister2025distributional} go through the same frozen rules.
The ladder keeps a few partisan gaps and leaves most other gaps failed, unresolved or unreadable, pointing toward narrower readings of each release's claims.
